\section{Neural Operators in One Page}\label{sec:fno-primer}

The Markov interlude showed that a commutative merge cannot carry the
boundary information the oracle depends on, so the learned summarizer
$g$ has to be a position-aware operator. This section records \emph{one
practical route} for obtaining such a $g$ (and, in the Markov
experiments, a readout $f$ baked into its training head): a Fourier
neural operator (FNO) backed by the equation-(6) neural-operator
architecture of \citet{Kovachki23NeuralOperator}. Neural operators give
us a parametric class with explicit approximation guarantees on compact
realized calls and a clean bridge to the paper's local-law budgets
(Proposition~\ref{prop:neural-operator-bridge}), which is why we use
them in the controlled benchmark. They are not, however, the only
route. Section~\ref{sec:manifesto-llm} instantiates $f$ and $g$ as
prompted LLM calls --- a different parametric family, optimized without
gradient-based operator training --- and the same local-law audit and
gap bound still apply because they are stated on $f$ and $g$, not on a
particular parameterization. The contribution here is the approximation-theoretic interface: an
operator family that can be applied across the realized leaf, merge,
and re-summary calls of one tree, with transfer moduli converting its
tolerance $\varepsilon$ into the budgets
$(\varepsilon_{\mathrm{leaf}}, \varepsilon_{\mathrm{merge}},
\varepsilon_{\mathrm{idemp}})$ of Section~\ref{sec:theorems}.

The architecture surface we use is the equation-(6) neural-operator form
from \citet{Kovachki23NeuralOperator}: a pointwise lift, a finite stack
of local-plus-kernel layers, and a pointwise projection,
\begin{equation}\label{eq:equation6-no}
\mathcal{G}_{\theta}
  =
  Q \circ
  \sigma_T(W_{T-1}+K_{T-1}+b_{T-1})
  \circ \cdots \circ
  \sigma_1(W_0+K_0+b_0)
  \circ P .
\end{equation}
We use this form only as a function-space interface. The FNO is one
implementation of the kernel layers, not an additional theorem premise.

A notation map relating Kovachki's operator-learning vocabulary to the
C-TreePO objects used here --- leaf calls, merge calls, on-range
re-summary calls, and the corresponding index domains --- is in
Appendix~\ref{app:operator-overlap} (Table~\ref{tab:ctreepo-neural-operator-notation}).
The one shift worth flagging upfront: our ``domain'' is not a PDE mesh
alone but the index set of the realized call currently being processed.

The first theoretical ingredient is \emph{discretization invariance} in
the sense of \citet[Section~2.3]{Kovachki23NeuralOperator}: one
parameter-sharing operator family can be evaluated on multiple
discretizations of the same underlying object, and refinement improves
the approximation rather than changing the target. This is the
manuscript's reusable requirement on the operator family: a single
learned operator can be shared across short leaves, longer merge inputs,
and re-summarized root states without changing the theorem surface.

The second ingredient is Section~9's approximation theorem. We separate
its two roles. First, Kovachki Lemma~22 finite-dimensionalizes a
continuous operator on a compact set: under the approximation-property
side conditions, $\mathcal{G}^\dagger$ can be written up to tolerance as
$G_{J'}\circ\phi\circ F_J$, where $F_J$ is a finite list of input
functionals, $\phi$ is a continuous finite-dimensional map, and $G_{J'}$
is a finite output expansion.\footnote{Paraphrased form of
\citet[Lemma~22]{Kovachki23NeuralOperator}: if $X,Y$ are Banach spaces
with the approximation property and $G:X\to Y$ is continuous, then for
every compact $K\subset X$ and $\varepsilon>0$ there exist
$J,J'\in\mathbb{N}$, continuous linear maps
$F_J:X\to\mathbb{R}^J$ and $G_{J'}:\mathbb{R}^{J'}\to Y$, and
$\phi\in C(\mathbb{R}^J;\mathbb{R}^{J'})$ such that
$\sup_{x\in K}\|G(x)-(G_{J'}\circ\phi\circ F_J)(x)\|_Y\le\varepsilon$,
with $F_J(x)=(w_1(x),\ldots,w_J(x))$ and
$G_{J'}(v)=\sum_{j=1}^{J'}v_j\beta_j$. Lemma~21 supplies the compact
union of finite-rank input images used to obtain the required uniform
continuity before choosing the finite-dimensional approximation.}
Appendix~\ref{app:kovachki-finite-dimensionalization} gives the exact proof
route we use for this finite-dimensionalization step, including the
uniform-continuity specialization that discharges Lemma~21's stability
premise when the target operator is uniformly continuous or Lipschitz on
$K$.

Second, the architecture theorem approximates the finite-dimensional
map $\phi$ by an equation-(6) neural operator. The paper uses only the
fragment relevant to C-TreePO: on a compact realized-call set and under
the usual Banach-space side conditions, a target operator can be
approximated to any desired tolerance.\footnote{Paraphrased
form of \citet[Theorem~11]{Kovachki23NeuralOperator}: under their
Assumptions~9--10, for a continuous target operator
$\mathcal{G}^\dagger:\mathcal{A}\to\mathcal{U}$, every compact
$K\subset\mathcal{A}$, and every tolerance $0<\varepsilon\le1$, there
exists an equation-(6)-style neural operator $\mathcal{G}$ with
$\sup_{a\in K}\|\mathcal{G}^\dagger(a)-\mathcal{G}(a)\|_{\mathcal{U}}
\le \varepsilon$.} When the downstream text uses expected risk rather
than uniform-on-compact control, we use the corresponding
Bochner-risk interface.\footnote{Paraphrased form of
\citet[Theorem~13]{Kovachki23NeuralOperator}: under the paper's
Banach/Hilbert and activation assumptions, if
$\mathcal{G}^\dagger$ is square-integrable with respect to a probability
measure $\mu$, then for every $0<\varepsilon\le1$ there is a neural
operator whose $L^2_{\mu}$ Bochner error is at most $\varepsilon$.}
Together these two statements give the ``arbitrary desired error
tolerance'' on compact realized-call sets that the next step consumes;
the equation-(6) architecture universal approximation itself is taken
as a named external input from Kovachki Theorem~11/13.

For C-TreePO the relevant compact sets are not all possible strings.
They are the compact realized leaf calls, compact realized merge calls,
and compact on-range re-summary calls induced by the tree we actually
run. That restriction matters. The approximation theorem is used only on
the support the pipeline visits, which is why the formal bridge works
with compact realized-call sets rather than an unrestricted all-strings
universe.

The third ingredient is the bridge from approximation error to local-law
budgets. Uniform operator approximation does \emph{not} automatically
imply the approximate local laws, because the latter are phrased in
terms of audited violation probabilities. The bridge of
Proposition~\ref{prop:neural-operator-bridge} therefore takes three
inputs explicitly: an ideal deterministic summarizer that is already
exact theorem-backed, a uniform approximation witness on the compact
realized leaf/merge/on-range call sets, and transfer assumptions that
convert approximation tolerance $\varepsilon$ into leaf, merge, and
idempotence budgets. This is the precise sense in which the
neural-operator layer is the \emph{realizability layer}, while C1--C3
remain the \emph{audit and certification layer}.
Because this route certifies deterministic operators, it enters the
general C-TreePO semantics through the point-mass specialization of the
stochastic summarizer interface. A randomized LLM summarizer is therefore
not certified by the neural-operator approximation theorem itself; it is
covered by the PMF-valued local-law theorem stack, or by an additional
randomized-operator approximation result if one is supplied.

The Markov experiments instantiate this interface with a 1D FNO: token
embeddings are mixed globally in Fourier space, pooled at leaves, and
then passed through a learned merge map. FNOs are useful here because
global frequency mixing makes boundary interactions available anywhere in
the window, which is exactly what the commutative-sketch baseline loses.
But the theorems below do not require Fourier structure specifically.
Any learned operator family satisfying the same realized-call
approximation interface and the same audited transfer conditions can
replace it unchanged.

The transformer statement we use is architecture-level. Following
\citet[Proposition~6]{Kovachki23NeuralOperator}, a finite
single-head, pre-normalized attention stack --- assembled from
attention, residual/local maps, and pointwise feed-forward maps ---
instantiates the equation-(6) architecture: single-head attention
agrees with a discretized kernel layer on the block's active
coordinates, transformer blocks are neural-operator layers, and a
finite encoder stack is an equation-(6) neural operator. This means a
transformer summarizer can be placed in the same operator-learning
pipeline as the FNO: external equation-(6) approximation bounds supply
the upstream realized-call tolerance, Proposition~\ref{prop:neural-operator-bridge}
translates that tolerance into local-law budgets, and the audit checks
the budgets actually realized by the deployed system. The inclusion
does not certify tokenizers, runtime serving details, or a particular
checkpoint; those remain deployment details outside the theorem
surface. The manifesto section therefore uses the transformer
inclusion as an architecture-level membership statement, and still uses
the audited realized-call bridge as the load-bearing correctness step.
